\documentclass{article}
\usepackage{graphicx}
\usepackage{amsmath}
\usepackage{amssymb}
\usepackage{amsfonts}
\usepackage{hyperref}
\usepackage{exercises}
\usepackage{indentfirst}
\usepackage{circledsteps}
\usepackage{amsthm}
\usepackage{csquotes}
\usepackage{yhmath}
\usepackage{comment}
\usepackage{placeins}
\usepackage{geometry}
\usepackage{tikz}
\newcommand*{\oldneg}{\mathord{\sim}}
\newcommand{\Z}{\mathbb{Z}}
\newcommand{\N}{\mathbb{N}}
\let\oldemptyset\emptyset

\theoremstyle{definition}
\newtheorem{definition}{Definition}

\setlength{\parindent}{.5in}
\geometry{margin=0.5in}
\setlength{\parskip}{12pt}
\title{Chapter 4 Notes}
\begin{document}
\date{}
\maketitle

\section*{Section 1}

\subsection*{Notes}
Note that when Smullyan says: \textquote{Thus, for example, the class $V$ is partially ordered under inclusion,} by inclusion he means 
\textit{set inclusion} (that is, being a subset), not \textit{set membership}.

Here is the proof of part (2) of Proposition 1.3. For suppose that some element $x$ had two predecessors $y$ and $z$. Then 
$y<x$ and $z<x$. If $y<z$, then $y<z<x$, contrary to the fact that $y$ is the predecessor of $x$. Similarly it cannot be that 
$z<y$. But, because $\le$ is a linear ordering, either $y\le z$ or $z\le y$. In either case, because neither $y<z$ or $z<y$ hold, $y=z$, thus completing 
the proof.

Note that, for a linear ordering $\le $ of $A$, all finite non-empty subclasses of $A$ must have a least element. We prove this by induction on the size 
of the subclass. The base case is obvious (because, by $O_1$, $x\le x$ for all $x$ in the field of $\le$). 
Now suppose the result holds for some $n$. Consider some subclass $S$ of $A$ of size $n+1$. Then 
for any element $x\in S$, suppose $x$ is not the least element. Then $S-\{x\}$, by the induction assumption, has a least element $a$. Because $x$ 
is not the least element, $a\le x$ (by $O_4$), meaning $a$ is the least element of $S$. Thus, for every element $x\in S$, either $x$ is the least element of 
$S$, or there is some other element least element $a$ of $S$. In any case, $S$ has a least element.

Here is an explicit proof of Proposition 1.7: We can prove by induction that, for a well-ordering $\le$, every non-empty finite subclass of the field of 
$\le$ has a greatest element. The base case, when the subclass has one element, is obvious (because, by $O_1$, $x\le x$ for all $x$ in the field of $\le$). 
Now suppose that, for some non-empty subclass of size $n$, it 
has a greatest element. Now consider some non-empty subclass $C$ of size $n+1$. By $O_5$, it has a least element $e$. Then $C-\{e\}$ is of size $n$, and thus 
by the induction assumption, has a greatest element $a$. But because $e$ is the least element of $C$, $e\le a$. Hence, for all elements $x$ of 
$C$, $x\le a$, meaning $a$ is the greatest element of $C$. This completes the induction.

Now having proven this result, suppose that $x$ was a limit element of a well ordering $\le$ and $L_<(x)$ had finitely many elements. Then, by the above result, 
$L_<(x)$ must have a greatest element, and thus, by Proposition 1.6, $x$ would not be a limit element, a contradiction.

\subsection*{Exercise 1.1}
Suppose that $\le$ is a partial ordering such that every non-empty subclass of the field of $\le$ has a least element. We need to show that, for every 
$x$ and $y$ in the field of $\le$, $x\le y\vee y\le x$. If $x$ and $y$ are both in the field of $\le$, then we know that the class 
$\{x,y\}$ has a least element. If $x$ is the least element, then $x\le y$, and if $y$ is the least element, then $y\le x$. Thus 
either $x\le y$ or $y\le x$, therefore completing the proof.

\subsection*{Exercise 1.2}
I do not think the premise of this problem makes sense. For suppose that a linear ordering $\le$ of $A$ was such that, for every proper lower section $L$ of 
$A$, there is a least element $x$ of $A$ not in $L$. In fact, we can then prove that $A$ has no proper lower sections. For suppose it has some proper lower 
section $L$. 
Then there would be a least element $x$ of $A$ not in $L$. Thus, for every element $y\in L$, $y\le x$. But, by the definition of a proper lower section, 
also $x<y$, which contradicts $O_2$. Hence $A$ can have no proper lower sections. Then obviously $\le$ cannot be a well-ordering of $A$. For if it was, then 
$A$ would have a least element $e$. But then $\{e\}$ would be a proper lower section of $A$, contrary to the fact that $A$ has no proper lower sections. Thus 
$\le$ cannot be a well-ordering of $A$.

Now, if the problem instead said: for every proper lower section $L$ of $A$, there is a least element $x$ of $A-L$, the situation would be different. Here the problem 
would not be contradictory, but the problem is not true. For instance, we can give the following counterexample. Consider the relation $x\supseteq y$ for all natural numbers $x$ and $y$. 
Clearly this is a linear ordering. However, there is no natural number that is greater than all others. Thus $\supseteq$ is not a 
well ordering of the natural numbers. However, it is the case that, for every proper lower section $L$ of $\omega$, there is a least element of $A-L$ (that is, a natural number $n$ of $A-L$ such 
that $n\supseteq x$ for all $x\in A-L$). For consider some proper lower section 
$L$ of $\omega$. Then it is a set of natural numbers, and hence has an element $n$ such that $n\subseteq x$ for all $x\in L$. Then $n-1$ is the least element of $A-L$, that is, 
$n-1\supseteq x$ for all $x\in A-L$. Thus, in fact, 
the problem does not hold for this particular linear ordering, meaning it is not true.

Now consider any linear ordering $\le$ for some class $A$. If we were given (in addition to the premises of the problem), that $A$ itself had a least element $a$, the desired 
result would follow. For then consider any non-empty subclass $C$ of $A$. If it was the entirety of $A$, then it would have a least element (by assumption). So suppose 
$C\subset A$. Then let $S$ be the class of all elements of $A$ that are less than every element in $C$. If $S$ is empty, then $a\in C$, in which case 
$a$ is the least element of $C$. Now suppose $S$ is not empty. Because $S$ is a proper lower section of 
$A$, it follows that $A-S$ has a least element $e$. But suppose $e\notin C$. Then $e$ would be an element that is $\le$ all elements in $C$ but not equal to any of them, meaning 
$e$ is less than all elements of $C$; thus, by the definition of 
$S$, $e\in S$, contrary to the fact that $e\in A-S$. Thus $e\in C$, meaning $e$ is the least element of $C$. Thus the proof is complete.

\subsection*{Exercise 1.3}
This problem had similar issues as in Exercise 1.2: assumption (c) of the problem implies that there can be no proper lower sections $L$ of $A$ such that 
$L$ has no greatest element.

So instead I am assuming that assumption (c) means: For any proper lower section $L$ of $A$, if $L$ has no greatest element, then there is a least 
element of $A-L$.

Suppose that $\le$ is a linear ordering of $A$.

(1): On the one hand, suppose that $\le$ is a well ordering of $A$. (a) obviously holds (by $O_5$). As for (b), consider any element $x\in A$ which is not the greatest element 
in $A$. If the set of all elements of $A$ that are $\ge x$ was empty, then $x$ would be the greatest element of $A$. Thus the set $S$ of all elements of $A$ that are $\ge x$ is non-empty. Because 
$\le$ is a well ordering, $S$ must contain a least element $e$. Then obviously $e$ is the successor of $x$. As for (c), consider any proper lower section $L$ of $A$ such that 
$L$ has no greatest element. Because it is a proper lower section, $A-L$ is non-empty, and hence, because $\le$ is a well ordering, contains a least element. This completes the first part of the 
proof.

(2): On the other hand, suppose that conditions (a-c) all hold. We must show that $\le$ is a well ordering for $A$. By (a), $A$ has some least element $a$. 
Now consider any non-empty subclass $C$ of $A$. If $C=A$, then it would have a least 
element (by assumption). So suppose $C\subset A$. Then let $S$ be the class of all elements of $A$ that are less than every element in $C$. If $S$ is empty, then $a\in A$, 
in which case $a$ is the least element of $C$. Now suppose $S$ is not empty. Clearly $S$ is a proper lower section of $A$. If $S$ had no greatest element, 
then, by (c), $A-S$ would contain a least element $e$ (and of course $e\in C$). On the other hand, suppose that $S$ did contain a greatest element $b$ (and $b$ cannot be the greatest element of 
$A$). Then, by (b), it has a successor $c$. Of course, $c$ cannot be in $S$, as that would contradict the fact that $b$ was the greatest element of $S$. 
First, $c$ must be in $C$. For suppose it weren't. Then there is some element $x\in C$ such that $x<c$ (because if this weren't true, then 
$c$ would be in $S$). Then $b<x<c$, contrary to the fact that $c$ is the successor of $b$. Thus $c\in C$. Second, $c$ is the least element of $C$. For suppose it weren't. Then there would be 
some element $y\in C$ such that $y<c$. But then $b<y<c$, contrary to the fact that $c$ is the successor of $b$. Thus $c$ is the least element of $C$. Thus every non-empty subclass of $A$ contains a least element, 
meaning $\le$ is a well ordering of $A$.

\subsection*{Exercise 1.4}
We prove that in a well ordering, all backward paths must be finite. What this amounts to saying is that, if $\le$ is a well-ordering, then 
there can be no denumerable backward paths. So suppose there was some denumerable backwards path $a_0,a_1,...,a_n,...$. Because this is a well-ordering, 
the class $\{a_0,a_1,...,a_n,...\}$ has a least element. Thus for some $m$, it must be the case that $a_m$ is the least element of 
$\{a_0,a_1,...,a_n,...\}$. Because the sequence is denumerable, there must be a term $a_{m+1}$ of the sequence. By the definition of a backward path, 
$a_{m+1}<a_m$. But also $a_m\le a_{m+1}$, which is a contradiction (by $O_2$, as this would imply $a_m=a_{m+1}$, but because 
$a_{m+1}<a_m$, $a_m\ne a_{m+1}$).

\subsection*{Exercise 1.5}
(a): On the one hand, we assume the first form of the AC. Now consider any non-empty set $A$ of pairwise disjoint non-empty sets. 
By the AC, there is some function $f$ such that $f(x)\in x$ for all elements $x\in A$. Then take $C$ to be the class of all $x$ such that 
$x=f(y)$ for some $y\in A$. Because $A$ is a set of pairwise disjoint sets, and $f(y)\in y$ for every $y\in A$, it follows that $C$ contains exactly one element of 
each element of $A$. 


(b): On the other hand, suppose that the second form of the AC holds. Now consider any non-empty set $A$ of non-empty sets. 
Note that for any set $x$, $\{x\}$ is a non-empty set of pairwise disjoint non-empty sets, and hence, there exists some set $C$ that consists of exactly one element 
of $x$. Denote this set $C$ for this set $x$ as $C_x$, and let $a_x$ be the one element of $C_x$. Then take the function $f$ to be the set of all ordered pairs
$\langle x,a_x\rangle$, which is clearly a function such that for each $x$ in $A$, $f(x)\in x$.

\subsection*{Exercise 1.6}
It suffices to show that this form is equivalent to the first form, and then, by Exercise 1.5, it follows that it is also equivalent to the form in that 
problem.

(a): On the one hand, suppose that the first form of the AC holds, and consider some non-empty set $A$ of non-empty sets. 
Then, by the AC, there is some function $f$ such that 
$f(x)\in x$ for all elements $x$ of $\mathcal{P}(A)-\{A\}$ (that is, the set of all proper subsets of $A$), that is, for all proper subsets $x$ of $A$, 
$f(x)\in x$.

(b): Now suppose that the second form (the one presented in this problem) 
of the AC holds, and consider some non-empty set $A$ of non-empty sets. Then $\cup A$ is a set, and hence, by the 
AC, there is some function $f$ such that $f(x)\in x$ for all proper subsets of $\cup A$. Then we take the function $g$ to be the set of all ordered pairs 
$\langle x,f(x)\rangle$ for all elements $x$ of $A$, which works, as $x$ is a subset of $\cup A$. 

\subsection*{Exercise 1.7}
We wish to prove that every infinite set can be put into a 1-1 correspondence with a proper subset of itself  (assuming the AC). So, consider any infinite set $A$. Then, by exercise 1.8,
it has a denumerable subset $D$, and there is a 1-1 function $f$ from the natural numbers to elements of $D$. Pick some element of $D$ ($f(0)$, say). Then 
$D-f(0)$ is still a denumerable subset. Now we can define a 1-1 function $g$ from $A$ to $A-\{f(0)\}$ (a proper subset of $A$) as follows: 
Assume $x\in A$. If $x\in D$, then $x=f(n)$ for some $n$, and we add the ordered pair $\langle f(n),f(n^+)\rangle$ to the function $g$. If $x\notin D$, 
then we add the ordered pair $\langle x,x\rangle$ to $g$. Thus this a 1-1 function $g$ from $A$ onto $A-\{f(0)\}$, therefore completing the proof. 

\subsection*{Exercise 1.8}
It can easily be shown by induction that every infinite set $A$ contains an infinite amount of subsets of size $n$ for every $n>0$. The base case, when $n=1$, holds, as, 
because $A$ is infinite, there is, for every element $x$ of $A$, the subset $\{x\}$ of $A$, and there are infinitely many of these. Now suppose the result 
holds for some $n$, that is, $A$ has an infinite amount of subsets of size $n$. For each of these subsets $S$ of size $n$, $A-S$ (which must also be infinite) has an 
infinite number of subsets of size 1 (by having proven the base case), and then, for each element $x\in A-S$, the subset $S\cup\{x\}$ is of size $n+1$, and because 
there are infinitely many of these (because $A-S$ contains infinitely many elements), the induction is complete.

Now consider any infinite set $A$. Then the class $C$ of all subsets of $A$ of size one 
is a set (because it is a subset of $\mathcal{P}(A)$). Then
$\{C\}$ is a set, meaning, by the AC, there is some function $f^0_0$ such that $f^0_0(C)\in C$, that is, $f^0_0(C)$ is a subset of $A$ 
of size one. Because 
$f^0_0(C)$ is a set, so is $\{f^0_0(C)\}$ (which is a non-empty set of non-empty sets), 
meaning, by the AC, there is some function $f^0_1$ such that $f^0_1(f^0_0(C))\in f^0_0(C)$, that is, 
$f^0_1(f^0_0(C))$ is an element of $A$ ($f^0_1(f^0_0(C))$ is an element of a subset of $A$ of size one). 
Then we can establish a denumerable sequence $a_0,a_1,...,a_n,...$ of elements of $A$ as follows: Take 
$a_0$ to be $f^0_1(f^0_0(C))$, and let $S_0=A$ and $C_0=C$. 
Then let $S_1=S_0-\{a_0\}$ and $C_1$ be the set of all subsets of $S_1$ of size one, 
and, in a similar manner as in how we obtained $f^0_1(f^0_0(C))$, we can define functions $f^1_1$ and $f^1_0$ using the AC such that $f^1_1(f^1_0(C_1))$
is an element of $S_1$ not equal to $a_0$ (because $S_1$ does not contain $a_1$). 
Then take $a_1$ to be $f^1_1(f^1_0(C_1))$. In general, if we have $a_i$ and $S_i$ for all numbers $i$ less than or equal to 
$n$, we take $S_{n+1}$ to be $S_n-\{a_n\}$, $C_{n+1}$ to be the set of all subsets of $S_{n+1}$ of size one, 
and we can define functions $f^{n+1}_1$ and $f^{n+1}_0$ using the AC such that 
$f^{n+1}_1(f^{n+1}_0(C_{n+1}))$ is an element of $S_{n+1}$ not equal to $a_i$ for all $i\le n$. Obviously, each $a_i$ is distinct from all others in the 
sequence. Thus the set of all $a_i$ is a denumerable subset of $A$.

\section*{Section 2}

\subsection*{Notes}
Note that $\omega$ is not superinductive (under the successor function). This is because it is not the case that $\omega$ is closed under chain unions. 
For instance, consider the subset
$C:\{1,2,...,n,...\}$ of $\omega$, which is obviously a chain. We show that it is not the case that $\cup C\in\omega$. For suppose $\cup C\in\omega$. Then, by definition 
of $\omega$, $\cup C$ would be some natural number $n$. Obviously $n\ne 0$. Now, $n^+\notin n$ (by Theorem 5.6 of Chapter 3). But $n^{++}$, which includes 
$n^+$, is an element of $C$; thus $n^+$ is an element of $\cup C$, a contradiction. Hence $\cup C\notin\omega$.

Here is an explanation of why $M$ is minimally superinductive under $g$ in Theorem 2.8. We know that $M$ is the intersection of all 
subsets of $S$ that are superinductive under $g$. Now, let $C$ be the set of all subsets of $S$ that are superinductive under $g$. 
Then $M=\cap C$. Because every element of $C$ is superinductive under $g$, every element of $C$ is closed under $g$, and every element of $g$ 
is closed under chain unions. Thus, by Lemma 2.7, $\cap C$, that is, $M$, is closed under $g$ and is closed under chain unions. Obviously 
$\emptyset\in M$ (because $\emptyset$ is an element of all superinductive sets, and hence of all elements of $C$). Lastly, we must show that $M$ is minimally 
superinductive under $g$. Suppose that some proper subset $A$ of $M$ was superinductive under $g$. Thus there must be some element $a\in M$
such that $a\notin A$. Now, clearly $A$ must be a subset of $S$. Thus $A$ is a subset of $S$ which is superinductive under $g$. But then $A\in C$, and because 
$a\in M$ ($a\in\cap C$), $a\in A$, a contradiction. Thus $M$ must be minimally superinductive under $g$ (and $M$ is a set because $M\subseteq S$). 

Here is why $\cup M\in S$ in the proof of Theorem 2.8. As remarked in the proof, we know that $\cup M\in M$. But $M$ is the intersection of all 
subsets of $S$ that are superinductive under $g$. But because $S$ itself is superinductive under $g$ (and is a subset of itself), 
$\cup M\in M$ implies that $\cup M\in S$. 

There are multiple things to be noted for the proof of Theorem 2.10. First, the definition of $S_b$ is not correct. For, given the current definition, we cannot 
definitively know that $S_b$ is closed under $g'$. For suppose that $x\in S_b$. Then either $x\ne\emptyset$ or $x=\emptyset$. In the first case, we know that 
$g'(x)=g(x)$, and because $x\subseteq g(x)$ (because $g$ is progressing), $x\subseteq g'(x)$. Thus, because $x$ is a set which includes $b$ (by the 
definition of $S_b$), $g'(x)$ also includes $b$, meaning $g'(x)\in S_b$. However, if $x=\emptyset$, then $g'(x)=b$. But there is no guarantee that 
$b\in S_b$: it very well could be the case that $b\notin b$. In fact, it very well could be that $S_b$ (defined in this way) only contains the empty set!

I will remark that this (incorrect) definition of $S_b$ does make $S_b$ closed under chain unions. For consider any chain $C$ of elements of $S_b$. Then 
consider $\cup C$. Each element of $C$ is either $\emptyset$ or a set $x$ such that $b\in x$. If $C$ only contains $\emptyset$, then $\cup C=\emptyset$, and 
thus $\cup C\in S_b$. Otherwise, $C$ contains some element $x$ such that $b\in x$. But then $\cup C$ contains all elements of $x$, and hence contains $b$, meaning 
$b\in\cup C$, and thus $\cup C\in S_b$.

The proof of Theorem 2.10 which I found is very similar to the one given in the book, just with a slightly different definition of $S_b$:

\textbf{Corrected Proof of Theorem 2.10}:\newline
Assume the hypothesis. As in the original proof, the case when $b=\emptyset$ is just Thoerem 2.8; so assume $b\ne\emptyset$. We let $S_b$ be the set of all 
elements $x$ of $S$ for which $b\subseteq x$ (and this \textit{does} include $b$), together with $\emptyset$. This is closed under chain unions for the following reason: 
Consider some chain $C$ of elements of $S_b$. Then consider $\cup C$. Each element of $C$ is either $\emptyset$ or a set that $b$ is a subset of. If $C$ only contains 
$\emptyset$, then $\cup C=\emptyset$, and thus $\cup C\in S_b$. Otherwise, $C$ contains some element $x$ such that $b\subseteq x$. But then $\cup C$ contains all 
of the elements of $x$, and thus also all of the elements of $b$, meaning $b\subseteq\cup C$, and thus $\cup C\in S_b$. Thus $S_b$ is closed under chain unions. 
Then we define $g'$ in the same way as in the book: $g'(\emptyset)=b$, and for any $x\in S_b$ other than $\emptyset$, take $g'(x)=g(x)$. Then $g'$ is progressing 
for the following reason: For consider any $x\in S_b$. If $x=\emptyset$, then obviously $x\subseteq g'(x)$. Otherwise, if $x\ne\emptyset$, then 
$g'(x)=g(x)$, and because $x\subseteq g(x)$ (because we assumed that $g$ was progressing), $x\subseteq g'(x)$. Thus $g'$ is progressing. Next, we need to show that 
$S_b$ is closed under $g'$. For suppose that $x\in S_b$. If $x=\emptyset$, then $g'(x)=b$, and because $b\in S_b$, $g'(x)\in b$. Now suppose 
$x\ne\emptyset$. Because $x\in S_b$ (and $x\ne\emptyset$), $b\subseteq x$. Thus, because $x\subseteq g'(x)$, $b\subseteq g'(x)$, meaning (by the definition 
of $S_b$), $g'(x)\in S_b$. Thus $S_b$ is closed under $g'$. Lastly, here is why $S_b$ is a set: $S_b-\{\emptyset\}$ is a set because it is a subset of 
$S$ (which is assumed to be a set), and thus $\cup\{(S_b-\{\emptyset\}),\{\emptyset\}\}$, that is, $S_b$, is also a set. Thus $S_b$ satisfies the criteria of 
Theorem 2.8, meaning there is some element $x$ of $S_b$ such that $g'(x)=x$. Now, it cannot be that $x=\emptyset$, or else 
$b=g'(x)=x=\emptyset$. Thus $x\ne\emptyset$, meaning (by the definition of $S_b$) $b\subseteq x$. 
Also, because $x\ne\emptyset$, $g'(x)=g(x)$, so that $g(x)=x$. Thus $b\subseteq x$ and $g(x)=x$, therefore completing the proof (and obviously 
$x\in S$).

It very well could be that the book simply contained a typo, and meant to define $S_b$ in the way it was defined in the correct proof.

\subsection*{Exercise 2.1}
Consider any set $S$:

(1) Suppose that every element of $S$ is closed under a function $g$. Now consider $\cap S$. To prove that $\cap S$ is also closed under $g$, we need to show that 
if $x\in\cap S$, then $g(x)\in\cap S$. So suppose $x\in\cap S$. Then $x$ is an element of every element of $S$. But every element of $S$ is closed under $g$. 
Thus $g(x)$ is an element of every element of $S$, meaning $g(x)\in\cap S$.

(2) Suppose that every element of $S$ is closed under chain unions. Now consider $\cap S$. To prove that $\cap S$ is also closed under chain unions, we need to show 
that for any chain $C$ of elements of $\cap S$, $\cup C\in A$. So consider some chain $C$ of elements of $\cap S$. Thus $C$ is also a chain of elements of every element of 
$S$. But, because every element of $S$ is closed under chain unions, $\cup C$ is an element of every element of $S$. Thus $\cup C\in\cap S$. 

\section*{Sections 3 and 4}

\subsection*{Notes}

In the proof of Theorem 3.1, there is no guarantee that $L$ is a set (because we only know for sure that $M$ is a class). 
Now, if we assumed that $M$ was a set, then $L$ certainly would also be a set. 
But this means that the proof simply is incorrect if we assume $M$ is only a class. This is a major problem with the proof. Luckily, the main theorem proved by it is that
the AC implies that every set can be well-ordered. Thus that theorem assumes we are dealing with sets, and so lessening the scope of Theorem 3.1 to only sets would still 
prove the main result.

In Theorem 3.1, here is why $M$ is linearly ordered under inclusion. In fact, every class is partially ordered under inclusion. Thus, by theorem 2.3, $O_4$ is also satisfied, 
meaning $M$ is linearly ordered under inclusion.

Note that $\emptyset$ can be well ordered under $\emptyset$.

In the discussion of the proof of Theorem 4.4, here is why $a\le b$ iff either $a=b$ or $a$ comes into $A$ before $b$. On the one hand, suppose that $a\le b$. Also suppose $a\ne b$. 
Then $F(a)\subset F(b)$. If $F(a)$ contained $b$, then $F(a)$ would be a lesser element of $N$ than $F(b)$ which contains $b$, a contradiction. Thus $F(a)$ must not contain $b$. On 
the other hand, suppose that either $a=b$ or $a$ comes into $A$ before $b$. If $a=b$, obviously $a\le b$. Now suppose $a$ comes into $A$ before $b$. 
Thus $F(a)$ does not contain $a$. Because $\le$ is a well ordering, either $a\le b$ or $b\le a$. If $a\le b$ did not hold, then we would have 
$b<a$, in which case $F(b)\subset F(a)$, meaning $F(a)$ did contain $b$, a contradiction. Thus $a\le b$.

\subsection*{Exercise 4.1}
(a): On the one hand, suppose that $a\le b$. Thus $F(a)\subseteq F(b)$. Now consider some element $x$ of $N$ that contains $b$. It cannot be that 
$x\subset F(a)$. For if that were true, then $F(a)\subseteq F(b)$ would be contradicted. Thsu $F(a)\subseteq x$, meaning $x$ contains $a$. On the other hand, 
suppose that every element of $N$ that contains $b$ also contains $a$. Now consider $F(a)$ and $F(b)$. If $F(b)\subset F(a)$, then, because $F(b)$ was assumed 
to contain $a$, $F(b)$ would be less than $F(a)$, contrary to the fact that $F(a)$ is the least element of $N$ that contains $a$. Thus 
$F(a)\subseteq F(b)$, meaning $a\le b$.

(b): This problem has a typo. It should say $x\in\cup N$, not $x\in N$: Suppose $x\in F(a)$ for some $x\in\cup N$ and $a\in\cup N$. Either 
$a\le x$ or $x\le a$. If $a<x$, then $F(a)\subset F(x)$. But it was assumed that $x\in F(a)$; thus $F(a)$ would be less than $F(x)$, a 
contradiction. Thus $x\le a$.

(c): We first show, for any elements $a$ and $b$ in $\cup N$, if $b$ is the successor of $a$, then $F(b)$ is the successor of $F(a)$. 
For suppose $b$ is the successor of $a$. Then $a<b$ and $F(a)\subset F(b)$, and for no $x\in\cup N$ is it the case that $a<x<b$. So suppose that 
for some $x\in\cup N$ $F(a)\subset F(x)\subset F(b)$. Then $a<x<b$, a contradiction. Thus $F(b)$ is the successor of $F(a)$.


Suppose $a\in\cup N$, and that $a$ is not the greatest element of $\cup N$. Now consider the successor $b$ of $a$ (and it must have some successor, as 
it is the least element of the class of all elements greater than $a$, which exists as this is a well ordering). For every $x\in \cup N-F(a)$, either 
$x\le b$ or $b\le x$. Suppose that $x<b$. Then $F(x)\subset F(b)$. Recall that, by the proof in the previous paragraph, $F(b)$ is the successor of 
$F(a)$. Now, 
either $F(x)\subseteq F(a)$ or $F(a)\subseteq F(x)$. If $F(a)\subset F(x)$, this would mean $F(a)\subset F(x)\subset F(b)$, contrary to the fact 
that $F(b)$ is the successor of $F(a)$. Thus $F(x)\subseteq F(a)$, meaning $x\in F(a)$. But $\cup N-F(a)$ contains no elements of $F(a)$. Thus $x$ 
is not in $\cup N-F(a)$. Hence for no element $x$ in $\cup N-F(a)$ is it the case that $F(x)\subset F(b)$, meaning $F(b)\subseteq F(x)$ for 
all all $x\in \cup N-F(a)$, and thus $b$ is the least element of $\cup N-F(a)$.

(d): This problem has a typo: It should say $a\in\cup N$, not $a\in N$. Suppose that $b$ is a successor element of $N$, and the predecessor of $b$ is 
$m$. We are to show that $b=F(a)$ for some $a\in\cup N$. 
Note that, because $N$ is slowly well ordered under $\subseteq$, $b$ contains exactly one more element $e$ than $m$. We can prove that 
$F(e)=b$. For suppose that there was some $y\in N$ that contained $e$ such that $y\subset b$. Now, either 
$y\subseteq m$ or $m\subseteq y$. But, because $e\notin m$, $m\subset y$. Thus $m\subset y\subset b$, contrary to the fact that 
$b$ is the successor of $m$. Hence $F(e)=b$. 

\subsection*{Exercise 4.2}
Suppose that for some set $S$, there is a choice function $f$ for $\mathcal{P}(S)$. We define the choice function $g$ for $S$ as follows: 
Consider some subset $C$ of $S$. The set $F$ of all elements 
$y$ such that $y=\{x\}$ for some $x\in C$ is a subset of $\mathcal{P}(S)$. Thus $f(F)\in F$, that is, $f(F)=\{x\}$ for 
some $x\in C$. Then we take $g(C)$ to be $x$. Thus $g$ is a choice function for $S$.

\subsection*{Exercise 4.3}
This I could not solve.

\end{document}
